\section{Results}
\label{sec:results}

\subsection{A 2.7M-parameter model against foundation models}
\label{sec:main}

Tables~\ref{tab:seizure} and~\ref{tab:events} report the four corpora
whose labels are paroxysmal events --- seizure detection on TUSZ and
CHB-MIT, epileptiform event and pattern classification on TUEV and IIIC
--- against seven supervised architectures and the foundation-model
baselines with their released weights and trained from scratch; the
remaining eight corpora are in Appendix~\ref{app:matrix} in the same
format. On seizure detection CroFreMo leads every baseline on the primary
metric by margins far outside the baselines' seed variance ($0.698$
vs.\ $0.545$ AUC-PR on TUSZ, $0.716$ vs.\ $0.627$ on CHB-MIT), and the
strongest baselines there are not the largest models. On TUEV it ties the
strongest pretrained foundation model and on IIIC it leads narrowly.
Across the eight remaining corpora it leads narrowly on TUEP, ties the
strongest supervised model on TUAB, and loses to a pretrained or a larger
from-scratch model on the two sleep-staging corpora, the two
cohort-diagnosis corpora, TUAR and Siena.

\begin{table}[t]
\centering
\scriptsize
\caption{Seizure detection. Mean $\pm$ std over three seeds; a superscript gives the seed count when fewer. Best per column in bold.}
\label{tab:seizure}
\resizebox{\textwidth}{!}{%
\begin{tabular}{lcccccc}
\toprule
\multirow{2}{*}{Method} & \multicolumn{3}{c}{TUSZ} & \multicolumn{3}{c}{CHB-MIT} \\
\cmidrule(lr){2-4} \cmidrule(lr){5-7}
 & Bal.\ Acc. & AUC-PR & AUROC & Bal.\ Acc. & AUC-PR & AUROC \\
\midrule
SPaRCNet & 0.5609 $\pm$ 0.0380 & 0.5020 $\pm$ 0.0213 & 0.8888 $\pm$ 0.0018 & 0.5127 $\pm$ 0.0079 & 0.1976 $\pm$ 0.0715 & 0.7995 $\pm$ 0.0374 \\
ContraWR & 0.5601 $\pm$ 0.0176 & 0.5043 $\pm$ 0.0175 & 0.8871 $\pm$ 0.0035 & 0.6362 $\pm$ 0.0168 & 0.5072 $\pm$ 0.0360 & 0.8905 $\pm$ 0.0239 \\
CNN-Transformer & 0.5389 $\pm$ 0.0548 & 0.3707 $\pm$ 0.1934 & 0.8455 $\pm$ 0.0420 & \textbf{0.6420 $\pm$ 0.0583} & 0.5057 $\pm$ 0.0741 & 0.9106 $\pm$ 0.0083 \\
FFCL & 0.5953 $\pm$ 0.0334 & 0.5449 $\pm$ 0.0296 & 0.8924 $\pm$ 0.0112 & 0.6163 $\pm$ 0.0146 & 0.5341 $\pm$ 0.0540 & 0.9124 $\pm$ 0.0107 \\
ST-Transformer & 0.5514 $\pm$ 0.0424 & 0.3589 $\pm$ 0.0314 & 0.8148 $\pm$ 0.0199 & 0.5551 $\pm$ 0.1040 & 0.2407 $\pm$ 0.1134 & 0.7208 $\pm$ 0.1037 \\
EEGNet & 0.5215 $\pm$ 0.0111 & 0.4921 $\pm$ 0.0302 & 0.8727 $\pm$ 0.0039 & 0.6016 $\pm$ 0.1455 & 0.1406 $\pm$ 0.1345 & 0.7442 $\pm$ 0.1100 \\
EEGConformer & 0.5216 $\pm$ 0.0374 & 0.3190 $\pm$ 0.0376 & 0.8186 $\pm$ 0.0283 & 0.5000 $\pm$ 0.0000 & 0.3054 $\pm$ 0.2530 & 0.5974 $\pm$ 0.1680 \\
\midrule
\multicolumn{7}{l}{\textit{Foundation models, released pretrained weights}} \\
BIOT & 0.5634 $\pm$ 0.0328 & 0.3379 $\pm$ 0.0379 & 0.8382 $\pm$ 0.0135 & 0.6190 $\pm$ 0.0447 & 0.5601 $\pm$ 0.0771 & 0.8738 $\pm$ 0.0036 \\
LaBraM-Base & 0.6139 $\pm$ 0.0378 & 0.4754 $\pm$ 0.0473 & 0.8687 $\pm$ 0.0252 & 0.5855 $\pm$ 0.0474 & 0.4519 $\pm$ 0.0969 & 0.8709 $\pm$ 0.0621 \\
CBraMod & 0.6347 $\pm$ 0.0342 & 0.4344 $\pm$ 0.0424 & 0.8520 $\pm$ 0.0050 & 0.5446 $\pm$ 0.0198 & 0.2332 $\pm$ 0.0978 & 0.7828 $\pm$ 0.0474 \\
EEGPT & 0.6587 $\pm$ 0.0619 & 0.5446 $\pm$ 0.0704 & 0.8821 $\pm$ 0.0122 & 0.5468 $\pm$ 0.0632 & 0.2811 $\pm$ 0.1325 & 0.8574 $\pm$ 0.0378 \\
TFM-Tokenizer & 0.5645 $\pm$ 0.0275 & 0.3474 $\pm$ 0.0185 & 0.8286 $\pm$ 0.0059 & 0.6344 $\pm$ 0.0163 & 0.6269 $\pm$ 0.0251 & \textbf{0.9291 $\pm$ 0.0055} \\
REVE-Base & 0.5523 $\pm$ 0.0372 & 0.5026 $\pm$ 0.0746 & 0.8885 $\pm$ 0.0158 & 0.6320 $\pm$ 0.0675 & 0.5802 $\pm$ 0.1956 & 0.8977 $\pm$ 0.0641 \\
\midrule
\multicolumn{7}{l}{\textit{Foundation models, trained from scratch}} \\
BIOT & 0.6156 $\pm$ 0.0645 & 0.3475 $\pm$ 0.0258 & 0.8256 $\pm$ 0.0166 & 0.6240 $\pm$ 0.0711 & 0.5661 $\pm$ 0.0722 & 0.8748 $\pm$ 0.0185 \\
LaBraM-Base & 0.5814 $\pm$ 0.0355 & 0.4752 $\pm$ 0.0477 & 0.8575 $\pm$ 0.0191 & 0.5854 $\pm$ 0.0649 & 0.3596 $\pm$ 0.0498 & 0.8226 $\pm$ 0.0298 \\
CBraMod & 0.5761 $\pm$ 0.0494 & 0.4818 $\pm$ 0.0522 & 0.8753 $\pm$ 0.0334 & 0.6018 $\pm$ 0.0883 & 0.3169 $\pm$ 0.2050 & 0.7802 $\pm$ 0.1333 \\
\midrule
CroFreMo (2.7M) & \textbf{0.6901$^{(1)}$} & \textbf{0.6979$^{(1)}$} & \textbf{0.9285$^{(1)}$} & 0.5308$^{(1)}$ & \textbf{0.7163$^{(1)}$} & 0.9288$^{(1)}$ \\
\bottomrule
\end{tabular}}
\end{table}

\begin{table}[t]
\centering
\scriptsize
\caption{Epileptiform event and pattern classification. Format as in Table~\ref{tab:seizure}.}
\label{tab:events}
\resizebox{\textwidth}{!}{%
\begin{tabular}{lcccccc}
\toprule
\multirow{2}{*}{Method} & \multicolumn{3}{c}{TUEV, 6-class} & \multicolumn{3}{c}{IIIC, 6-class} \\
\cmidrule(lr){2-4} \cmidrule(lr){5-7}
 & Bal.\ Acc. & Cohen's $\kappa$ & Weighted F1 & Bal.\ Acc. & Cohen's $\kappa$ & Weighted F1 \\
\midrule
SPaRCNet & 0.4654 $\pm$ 0.0346 & 0.4885 $\pm$ 0.0279 & 0.7354 $\pm$ 0.0146 & 0.4767 $\pm$ 0.0103 & 0.3719 $\pm$ 0.0118 & 0.4791 $\pm$ 0.0113 \\
ContraWR & 0.4195 $\pm$ 0.1258 & 0.4076 $\pm$ 0.0984 & 0.6962 $\pm$ 0.0455 & 0.5121 $\pm$ 0.0072 & 0.4135 $\pm$ 0.0053 & 0.5149 $\pm$ 0.0047 \\
CNN-Transformer & 0.3679 $\pm$ 0.1026 & 0.4098 $\pm$ 0.1165 & 0.6916 $\pm$ 0.0744 & 0.5148 $\pm$ 0.0084 & 0.4162 $\pm$ 0.0095 & 0.5131 $\pm$ 0.0143 \\
FFCL & 0.4376 $\pm$ 0.0687 & 0.4776 $\pm$ 0.0322 & 0.7366 $\pm$ 0.0169 & 0.5200 $\pm$ 0.0176 & 0.4231 $\pm$ 0.0197 & 0.5240 $\pm$ 0.0167 \\
ST-Transformer & 0.4731 $\pm$ 0.0199 & 0.5006 $\pm$ 0.0117 & 0.7411 $\pm$ 0.0087 & 0.4065 $\pm$ 0.0025 & 0.2961 $\pm$ 0.0021 & 0.4084 $\pm$ 0.0081 \\
EEGNet & 0.4152 $\pm$ 0.0575 & 0.4788 $\pm$ 0.0908 & 0.7389 $\pm$ 0.0367 & -- & -- & -- \\
EEGConformer & 0.4212 $\pm$ 0.0454 & 0.4548 $\pm$ 0.0041 & 0.7166 $\pm$ 0.0059 & -- & -- & -- \\
\midrule
\multicolumn{7}{l}{\textit{Foundation models, released pretrained weights}} \\
BIOT & 0.4894 $\pm$ 0.0302 & 0.5304 $\pm$ 0.0277 & 0.7532 $\pm$ 0.0158 & 0.4650 $\pm$ 0.0154 & 0.3605 $\pm$ 0.0163 & 0.4665 $\pm$ 0.0128 \\
LaBraM-Base & 0.6305 $\pm$ 0.0306 & 0.5984 $\pm$ 0.0031 & 0.8086 $\pm$ 0.0048 & 0.4830 $\pm$ 0.0177 & 0.3834 $\pm$ 0.0204 & 0.4852 $\pm$ 0.0189 \\
CBraMod & 0.6029 $\pm$ 0.0108 & 0.6449 $\pm$ 0.0217 & 0.8130 $\pm$ 0.0117 & 0.4161 $\pm$ 0.0151 & 0.3133 $\pm$ 0.0191 & 0.4178 $\pm$ 0.0179 \\
EEGPT & 0.5131 $\pm$ 0.0750 & 0.5736 $\pm$ 0.0396 & 0.7766 $\pm$ 0.0200 & 0.4965 $\pm$ 0.0179 & 0.3966 $\pm$ 0.0195 & 0.4996 $\pm$ 0.0138 \\
TFM-Tokenizer & 0.5959 $\pm$ 0.0359 & 0.6519 $\pm$ 0.0209 & 0.8171 $\pm$ 0.0105 & 0.5058 $\pm$ 0.0094 & 0.4060 $\pm$ 0.0090 & 0.5080 $\pm$ 0.0066 \\
REVE-Base & \textbf{0.6477 $\pm$ 0.0183} & \textbf{0.6846 $\pm$ 0.0389} & \textbf{0.8378 $\pm$ 0.0187} & 0.5290 $\pm$ 0.0032 & 0.4363 $\pm$ 0.0030 & 0.5338 $\pm$ 0.0025 \\
\midrule
\multicolumn{7}{l}{\textit{Foundation models, trained from scratch}} \\
BIOT & 0.5351 $\pm$ 0.0625 & 0.5324 $\pm$ 0.0250 & 0.7612 $\pm$ 0.0126 & 0.4601 $\pm$ 0.0137 & 0.3557 $\pm$ 0.0144 & 0.4618 $\pm$ 0.0151 \\
LaBraM-Base & 0.3724 $\pm$ 0.0853 & 0.3717 $\pm$ 0.0089 & 0.6928 $\pm$ 0.0133 & 0.5047 $\pm$ 0.0045 & 0.4057 $\pm$ 0.0067 & 0.5038 $\pm$ 0.0059 \\
CBraMod & 0.5744 $\pm$ 0.0105 & 0.5638 $\pm$ 0.0236 & 0.7671 $\pm$ 0.0041 & 0.4892 $\pm$ 0.0050 & 0.3935 $\pm$ 0.0058 & 0.4975 $\pm$ 0.0043 \\
EEGPT & -- & -- & -- & 0.4717 $\pm$ 0.0214 & 0.3680 $\pm$ 0.0236 & 0.4762 $\pm$ 0.0203 \\
TFM-Tokenizer & -- & -- & -- & 0.3848 $\pm$ 0.0064 & 0.2627 $\pm$ 0.0081 & 0.3806 $\pm$ 0.0094 \\
\midrule
CroFreMo (2.7M) & 0.5831$^{(1)}$ & 0.6804$^{(1)}$ & 0.8338$^{(1)}$ & \textbf{0.5405$^{(1)}$} & \textbf{0.4468$^{(1)}$} & \textbf{0.5497$^{(1)}$} \\
\bottomrule
\end{tabular}}
\end{table}

\subsection{What the seizure margin is made of}
\label{sec:attribution}

Tables~\ref{tab:abl-seizure} and~\ref{tab:abl-events} are the experiment the encoder was built for.
The axis-free model is trained with its tokenizer producing either the
duplex grid of Section~\ref{sec:tokenizer} or waveform tokens only, with
every other setting identical, at the adopted width and at one width
larger. Compared with the waveform-only tokenizer, the duplex tokenizer improves
CHB-MIT by $0.136$ AUC-PR, TUEV by $0.107$ $\kappa$, TUSZ by $0.063$ AUC-PR
and IIIC by $0.027$ $\kappa$ at $d{=}192$. At $d{=}256$ the four effects keep their sign; three
are of similar size and the CHB-MIT effect drops to $0.025$, so the size
of the effect is sensitive to capacity and optimization state, not only
to the corpus. For scale, the seed standard deviations of the same
recipe on the same corpora, measured on an earlier design with a
dedicated frequency axis (Appendix~\ref{app:triaxial}), are $0.055$, $0.036$, $0.035$ and $0.009$; they
are a reference for the recipe, not a substitute for this model's own
variance. The two margins by which the model leads the
foundation-model field are therefore not properties of the band
decomposition or of the folded attention alone: without the interaction
rows the same encoder falls to $0.635$ on TUSZ and $0.580$ on CHB-MIT,
and on CHB-MIT below TFM-Tokenizer.

Table~\ref{tab:perclass} breaks the TUEV effect down by class. Every
class improves, and the gain is largest on the two periodic-discharge
classes (GPED F1 $0.59 \to 0.83$, PLED $0.50 \to 0.58$), whose defining
morphology is fast activity locked to a slow rhythm; the rare
spike-and-sharp-wave class moves from near-zero recall to $0.16$, and
background, artifact and eye movement change little. The confusion that
remains is between GPED and PLED and between artifact and background.

The waveform-only comparison changes more than one thing at once ---
the waveform-only grid has half the rows and exposes neither an
amplitude envelope nor a phase feature --- so Table~\ref{tab:control}
adds the matched control that separates the analytic features from the
cross-frequency alignment. The extra rows carry, in turn, waveform
tokens, each band's own analytic features (its amplitude feature rotated
by its \emph{own} unit phase, so that no information from any other band
enters the row), or the interaction tokens of Eq.~(\ref{eq:rotation});
rows, gates, encoder, recipe and seed are otherwise identical, and we
verified on the frontend that under the own-phase variant perturbing the
phase of the slower bands leaves the faster bands' rows exactly
unchanged. Across the corpora with results, the analytic-only rows
recover at most a fraction of the duplex gain and on TUSZ fall below the
waveform-only grid; the cross-band alignment accounts for most of the
effect. This is the control that licenses the phrase ``cross-frequency
coupling as token content'' rather than ``an analytic-signal front end''.
The remaining controls of a fuller design --- amplitude-only rows, a
surrogate that scrambles which slow band drives which fast band, and a
learned fusion of the same ingredients --- are not in this version.

\begin{table}[t]
\centering
\scriptsize
\caption{Duplex vs.\ waveform-only tokenizer in the same encoder, seizure detection. Single seeds.}
\label{tab:abl-seizure}
\resizebox{\textwidth}{!}{%
\begin{tabular}{lcccccc}
\toprule
\multirow{2}{*}{Tokenizer, width} & \multicolumn{3}{c}{TUSZ} & \multicolumn{3}{c}{CHB-MIT} \\
\cmidrule(lr){2-4} \cmidrule(lr){5-7}
 & Bal.\ Acc. & AUC-PR & AUROC & Bal.\ Acc. & AUC-PR & AUROC \\
\midrule
waveform-only, $d{=}192$ & 0.6395$^{(1)}$ & 0.6347$^{(1)}$ & 0.9238$^{(1)}$ & \textbf{0.5670$^{(1)}$} & 0.5797$^{(1)}$ & 0.9079$^{(1)}$ \\
duplex, $d{=}192$ (final) & \textbf{0.6901$^{(1)}$} & \textbf{0.6979$^{(1)}$} & 0.9285$^{(1)}$ & 0.5308$^{(1)}$ & \textbf{0.7163$^{(1)}$} & \textbf{0.9288$^{(1)}$} \\
waveform-only, $d{=}256$ & 0.5728$^{(1)}$ & 0.5405$^{(1)}$ & 0.9144$^{(1)}$ & 0.5199$^{(1)}$ & 0.5826$^{(1)}$ & 0.9035$^{(1)}$ \\
duplex, $d{=}256$ & 0.5991$^{(1)}$ & 0.6311$^{(1)}$ & \textbf{0.9299$^{(1)}$} & 0.5399$^{(1)}$ & 0.6077$^{(1)}$ & 0.9120$^{(1)}$ \\
\bottomrule
\end{tabular}}
\end{table}

\begin{table}[t]
\centering
\scriptsize
\caption{Duplex vs.\ waveform-only tokenizer in the same encoder, event classification. Single seeds.}
\label{tab:abl-events}
\resizebox{\textwidth}{!}{%
\begin{tabular}{lcccccc}
\toprule
\multirow{2}{*}{Tokenizer, width} & \multicolumn{3}{c}{TUEV, 6-class} & \multicolumn{3}{c}{IIIC, 6-class} \\
\cmidrule(lr){2-4} \cmidrule(lr){5-7}
 & Bal.\ Acc. & Cohen's $\kappa$ & Weighted F1 & Bal.\ Acc. & Cohen's $\kappa$ & Weighted F1 \\
\midrule
waveform-only, $d{=}192$ & 0.4786$^{(1)}$ & 0.5730$^{(1)}$ & 0.7754$^{(1)}$ & 0.5159$^{(1)}$ & 0.4203$^{(1)}$ & 0.5150$^{(1)}$ \\
duplex, $d{=}192$ (final) & \textbf{0.5831$^{(1)}$} & \textbf{0.6804$^{(1)}$} & \textbf{0.8338$^{(1)}$} & 0.5405$^{(1)}$ & 0.4468$^{(1)}$ & 0.5497$^{(1)}$ \\
waveform-only, $d{=}256$ & 0.5660$^{(1)}$ & 0.6009$^{(1)}$ & 0.7918$^{(1)}$ & 0.5133$^{(1)}$ & 0.4179$^{(1)}$ & 0.5195$^{(1)}$ \\
duplex, $d{=}256$ & 0.5224$^{(1)}$ & 0.6792$^{(1)}$ & 0.8259$^{(1)}$ & \textbf{0.5550$^{(1)}$} & \textbf{0.4651$^{(1)}$} & \textbf{0.5558$^{(1)}$} \\
\bottomrule
\end{tabular}}
\end{table}

\begin{table}[t]
\centering
\scriptsize
\caption{Matched control: the same encoder, rows, gates and recipe, with the extra rows carrying waveform tokens, each band's own analytic features (amplitude and own phase, no cross-band alignment), or the interaction tokens. Primary metric; three seeds where $\pm$ is given, otherwise the seed count is superscripted.}
\label{tab:control}
\resizebox{\textwidth}{!}{%
\begin{tabular}{lcccc}
\toprule
Extra rows carry & TUEV (Cohen's $\kappa$) & TUSZ (AUC-PR) & CHB-MIT (AUC-PR) & IIIC (Cohen's $\kappa$) \\
\midrule
waveform-only & 0.5543$^{(2)}$ & 0.6346 $\pm$ 0.0093 & 0.6333$^{(2)}$ & 0.4336 $\pm$ 0.0127 \\
+ per-band analytic features, no cross-band alignment & 0.5901 $\pm$ 0.0139 & 0.6106 $\pm$ 0.0317 & -- & 0.4298$^{(1)}$ \\
+ interaction tokens (duplex, final) & \textbf{0.6546$^{(2)}$} & \textbf{0.6800 $\pm$ 0.0476} & \textbf{0.7163$^{(1)}$} & \textbf{0.4468$^{(1)}$} \\
\bottomrule
\end{tabular}}
\end{table}

\begin{table}[t]
\centering
\scriptsize
\caption{TUEV per class, CroFreMo with the waveform-only vs.\ the duplex
tokenizer (single seed each, 29,421 test windows): test-set count and
precision / recall / F1. Macro-F1 $0.48 \to 0.58$.}
\label{tab:perclass}
\begin{tabular}{lrccc}
\toprule
Class & $n$ & waveform-only P / R / F1 & duplex P / R / F1 & $\Delta$F1 \\
\midrule
SPSW & 567 & 0.87 / 0.04 / 0.07 & 0.29 / 0.16 / 0.20 & $+0.13$ \\
GPED & 4,677 & 0.61 / 0.56 / 0.59 & 0.82 / 0.83 / 0.83 & $+0.24$ \\
PLED & 1,998 & 0.48 / 0.53 / 0.50 & 0.63 / 0.54 / 0.58 & $+0.08$ \\
EYEM & 329 & 0.94 / 0.23 / 0.37 & 0.44 / 0.42 / 0.43 & $+0.06$ \\
ARTF & 2,204 & 0.35 / 0.62 / 0.44 & 0.42 / 0.65 / 0.51 & $+0.07$ \\
BCKG & 19,646 & 0.93 / 0.89 / 0.91 & 0.94 / 0.91 / 0.92 & $+0.01$ \\
\bottomrule
\end{tabular}
\end{table}


\subsection{The same token in other encoders}
\label{sec:content-structure}

Table~\ref{tab:content-structure} collects the tokenizer's effect across
four encoders. The interventions are not the same experiment, and the
table says which is which: for CroFreMo it is the duplex tokenizer against
a waveform-only one; for CBraMod it is interaction rows against matched
waveform rows added to the native tokenizer; for LaBraM and REVE it is the
whole CroFreMo tokenizer, band decomposition included, against the host's
own. CBraMod's encoder, which keeps its own rFFT magnitude branch, gains from
the interaction rows relative to matched waveform rows on TUEV ($0.045$),
TUEP ($0.057$), IIIC ($0.015$) and CHB-MIT ($0.124$), and loses $0.092$ on
TUSZ. LaBraM and REVE, trained from scratch with no spectral pathway,
gain $0.150$ and $0.241$ on TUEV when the tokenizer replaces their patch
embedding. CroFreMo gains on all four corpora measured.

\begin{table}[t]
\centering
\scriptsize
\caption{Tokenizer benefit across four encoders. The interventions differ
and the numbers are not one quantity: duplex vs.\ waveform-only tokenizer
(CroFreMo); interaction rows vs.\ matched waveform rows added to the
native tokenizer (CBraMod); the whole CroFreMo tokenizer vs.\ the native
tokenizer (LaBraM, REVE). Single seeds.}
\label{tab:content-structure}
\resizebox{\textwidth}{!}{%
\begin{tabular}{lllccccc}
\toprule
Encoder & Intervention & Control & TUEV & TUSZ & CHB-MIT & IIIC & TUEP \\
\midrule
CroFreMo & duplex tokenizer & waveform-only tokenizer & $+0.107$ & $+0.063$ & $+0.136$ & $+0.027$ & --- \\
CBraMod (scratch) & + interaction rows & + waveform rows & $+0.045$ & $-0.092$ & $+0.124$ & $+0.015$ & $+0.057$ \\
LaBraM (scratch) & CroFreMo tokenizer & native tokenizer & $+0.150$ & --- & $+0.017$ & $+0.014$ & --- \\
REVE (scratch) & CroFreMo tokenizer & native tokenizer & $+0.241$ & --- & --- & $+0.078$ & --- \\
\bottomrule
\end{tabular}}
\end{table}

Read within each intervention, the pattern is this. With its own
tokenizer replaced, LaBraM and REVE gain most on TUEV; that gain includes
the band decomposition and cannot yet be split between it and the
interaction rows, which would need the CroFreMo waveform-only tokenizer
in the same hosts. With the native tokenizer kept and rows added, CBraMod
gains on the three event-classification corpora and on CHB-MIT, where it
trains poorly from scratch and the interaction rows give it a foothold,
and loses on TUSZ. Whether an encoder with a spectral branch computes for
itself what the interaction rows supply is a question for probing or
intervention experiments; the table records that the benefit of explicit
cross-frequency content varies with the encoder and the corpus, not why.

\subsection{Design controls}
\label{sec:design}

Table~\ref{tab:design} reports the variants that were tried on the way
to the final model, all single seeds. Folding the band rows
into spatial attention rather than leaving them as unrelated channels is
worth $0.09$ on both seizure corpora and costs $0.02$--$0.04$ on the
smaller corpora, where the longer attention span overfits. Raising the
width from 192 to 256 recovers those small-corpus cells and loses
$0.07$ and $0.11$ on TUSZ and CHB-MIT. An explicit coupling-strength
feature --- the magnitudes $|Z_{ij}|$ projected onto the fused rows ---
raises TUEV to $0.721$ and Siena to $0.398$ and collapses CHB-MIT to
$0.513$. No variant dominates; the adopted one is the one that keeps the
seizure margin, and every alternative is reported. Two controls run on an earlier design with a dedicated frequency axis,
reported in Appendix~\ref{app:triaxial}, locate where a small model's
seizure-detection margin comes from: collapsing the learned filterbank to
one broadband band costs $0.11$ AUC-PR on TUSZ and $0.42$ on CHB-MIT, and
replacing factorized attention by a flat encoder costs $0.22$ and $0.53$.

\begin{table}[t]
\centering
\scriptsize
\caption{Design variants of the final model, primary metric per corpus, single seeds. ``fold'' = band rows folded into spatial attention; ``strength'' = explicit $|Z|$ feature.}
\label{tab:design}
\resizebox{\textwidth}{!}{%
\begin{tabular}{lcccccccc}
\toprule
Variant & TUSZ (AUC-PR) & CHB-MIT (AUC-PR) & TUEV (Cohen's $\kappa$) & IIIC (Cohen's $\kappa$) & TUAR (Cohen's $\kappa$) & ADFD (Bal.\ Acc.) & Siena (AUC-PR) & Sleep-EDF (Cohen's $\kappa$) \\
\midrule
no fold, $d{=}192$ & 0.6115$^{(1)}$ & 0.6302$^{(1)}$ & 0.6896$^{(1)}$ & 0.4577$^{(1)}$ & 0.6237$^{(1)}$ & \textbf{0.4982$^{(1)}$} & 0.2393$^{(1)}$ & 0.6099$^{(1)}$ \\
no fold, $d{=}192$, strength & 0.6429$^{(1)}$ & 0.5133$^{(1)}$ & \textbf{0.7207$^{(1)}$} & 0.4475$^{(1)}$ & 0.6261$^{(1)}$ & 0.4667$^{(1)}$ & \textbf{0.3982$^{(1)}$} & 0.6144$^{(1)}$ \\
fold, $d{=}192$ (final) & \textbf{0.6979$^{(1)}$} & \textbf{0.7163$^{(1)}$} & 0.6804$^{(1)}$ & 0.4468$^{(1)}$ & 0.5997$^{(1)}$ & 0.4608$^{(1)}$ & 0.2009$^{(1)}$ & 0.6393$^{(1)}$ \\
fold, $d{=}256$ & 0.6311$^{(1)}$ & 0.6077$^{(1)}$ & 0.6792$^{(1)}$ & \textbf{0.4651$^{(1)}$} & \textbf{0.6567$^{(1)}$} & 0.4540$^{(1)}$ & 0.2691$^{(1)}$ & \textbf{0.6529$^{(1)}$} \\
fold, strength, $d{=}192$ & 0.6711$^{(1)}$ & 0.6709$^{(1)}$ & 0.6606$^{(1)}$ & 0.4268$^{(1)}$ & 0.5452$^{(1)}$ & 0.3610$^{(1)}$ & 0.2800$^{(1)}$ & 0.6325$^{(1)}$ \\
\bottomrule
\end{tabular}}
\end{table}

\subsection{The token inside published encoders}
\label{sec:transplant}

Table~\ref{tab:transplant} adds the tokenizer's rows to CBraMod as
described in Section~\ref{sec:encoder}. Against the matched control ---
the same eight extra rows carrying waveform tokens --- the interaction
rows are better on TUEV, TUEP, IIIC and CHB-MIT and worse on TUSZ, so
the effect is attributable to what the rows carry rather than to the
extra channels. Against CBraMod's own tokenizer the transplant is a clear
gain on TUEV ($+0.077$, three standard deviations of the native model) and
on CHB-MIT ($+0.148$, where the native model is unstable across seeds), a
tie on TUEP and IIIC, and a loss on TUSZ.
Table~\ref{tab:hosts} widens the test to two further encoders, with
the tokenizer now \emph{replacing} the host's patch embedding (every
electrode--row pair becomes one host channel, rows are pooled after the
encoder, the head is unchanged). Trained from scratch, LaBraM gains
$0.150$ $\kappa$ on TUEV and REVE $0.241$; on IIIC the gains are
$0.014$ and $0.078$; on CHB-MIT LaBraM gains $0.017$, within its own
seed variance. The three hosts span 4M to 69M parameters and three
encoder designs, and the ordering of the TUEV gains follows how far each
host is from what it can learn from $10^5$ labelled windows on its own.
Adding the same rows to CBraMod's \emph{released} checkpoint is worse
than the checkpoint alone on both corpora tried (appendix): a pretrained
encoder is tuned to its own token distribution, and inserting eight new
channels per electrode disturbs it more than the coupling content
repays. We draw the conclusion the data support --- the token transfers
into encoders trained from scratch, its effect is attributable to the
content of the added rows, and it pays off on epileptiform event
classification --- and not the broad one about released checkpoints.

\begin{table}[t]
\centering
\scriptsize
\caption{CroFreMo's rows added to CBraMod's encoder, trained from scratch: the same eight extra rows per electrode carrying waveform tokens (control) or interaction tokens. Native: three seeds; transplants: single seeds.}
\label{tab:transplant}
\resizebox{\textwidth}{!}{%
\begin{tabular}{llccccc}
\toprule
Encoder & Tokenizer & TUEV (Cohen's $\kappa$) & TUEP (AUROC) & IIIC (Cohen's $\kappa$) & CHB-MIT (AUC-PR) & TUSZ (AUC-PR) \\
\midrule
CBraMod (from scratch) & native & 0.5638 $\pm$ 0.0236 & 0.6422 $\pm$ 0.0339 & 0.3935 $\pm$ 0.0058 & 0.3169 $\pm$ 0.2050 & \textbf{0.4818 $\pm$ 0.0522} \\
 & native + 8 waveform rows & 0.5956$^{(1)}$ & 0.5929$^{(1)}$ & 0.3815$^{(1)}$ & 0.3406$^{(1)}$ & 0.4762$^{(1)}$ \\
 & native + 8 interaction rows & \textbf{0.6407$^{(1)}$} & \textbf{0.6496$^{(1)}$} & \textbf{0.3962$^{(1)}$} & \textbf{0.4650$^{(1)}$} & 0.3839$^{(1)}$ \\
\bottomrule
\end{tabular}}
\end{table}

\begin{table}[t]
\centering
\scriptsize
\caption{CroFreMo's tokenizer in place of the host's own, hosts trained from scratch (every electrode--row pair becomes one host channel; rows pooled after the encoder; head unchanged). Native rows: three seeds except REVE (single); transplants: single seeds.}
\label{tab:hosts}
\resizebox{\textwidth}{!}{%
\begin{tabular}{llcccc}
\toprule
Encoder & Tokenizer & TUEV (Cohen's $\kappa$) & IIIC (Cohen's $\kappa$) & CHB-MIT (AUC-PR) & TUSZ (AUC-PR) \\
\midrule
CBraMod (4M) & native & 0.5638 $\pm$ 0.0236 & 0.3935 $\pm$ 0.0058 & 0.3169 $\pm$ 0.2050 & \textbf{0.4818 $\pm$ 0.0522} \\
 & CroFreMo tokenizer & \textbf{0.6322$^{(1)}$} & -- & -- & 0.4747$^{(1)}$ \\
\midrule
LaBraM (5.9M) & native & 0.3717 $\pm$ 0.0089 & 0.4057 $\pm$ 0.0067 & 0.3596 $\pm$ 0.0498 & 0.4752 $\pm$ 0.0477 \\
 & CroFreMo tokenizer & \textbf{0.5220$^{(1)}$} & \textbf{0.4202$^{(1)}$} & \textbf{0.3771$^{(1)}$} & -- \\
\midrule
REVE (69M) & native & 0.3033$^{(1)}$ & 0.2989$^{(1)}$ & -- & -- \\
 & CroFreMo tokenizer & \textbf{0.5437$^{(1)}$} & \textbf{0.3774$^{(1)}$} & -- & -- \\
\bottomrule
\end{tabular}}
\end{table}

\subsection{Pretraining}
\label{sec:pretrain-results}

CroFreMo is trained from scratch throughout. A pretraining study on an
earlier design with a dedicated frequency axis --- gains where labels are
scarce or noisy, losses where they are plentiful --- is reported as an
exploratory result in Appendix~\ref{app:pretrain}; we do not build on it.
